\eabstract{
This study aims to elucidate the long-term citation dynamics and structural characteristics of highly cited papers in the field of biomedical sciences. In conventional bibliometrics, short-term citation indicators such as the 3-year and 5-year impact factors have been predominantly used; however, long-term observation is essential for accurately capturing the true impact of scholarly articles. In this study, analysis was conducted from multiple perspectives, including citation lifespan, journal structure, and research topic structure, rather than relying solely on citation counts.
The dataset used was the PubMed Annual Baseline (2023), covering reference information from approximately 30 million papers published between 1946 and 2020. Based on this dataset, a “cumulative generation model” was constructed at five-year intervals, and the top 1\% of papers in terms of total citations within each cumulative generation were defined as “highly cited papers.” In addition, four groups were established to enable comparisons between highly cited papers and other papers.
At the individual paper level, the analysis revealed that even among highly cited papers, some of the most highly cited papers exhibited a long-lived pattern in which they did not reach obsolescence (citation decay) even several decades after publication. In contrast, among highly cited papers for which citation decline was observed, a lifespan structure was identified consisting of four years from publication to the citation peak and 12 years from publication to citation decay, which was consistent with existing models in the biological sciences.
At the journal level, it was shown that journals publishing many highly cited papers are not necessarily those with high impact factors. However, analysis using Kullback–Leibler divergence demonstrated that highly cited papers tend to exhibit a stronger concentration (preference) toward specific journals compared to other groups.
At the research topic level, entropy analysis based on MeSH descriptors revealed that highly cited papers are more concentrated in their research topics than other groups, and the topics of papers citing them are also similarly concentrated. At the same time, a divergence was observed between the topics of the cited papers and those of the citing papers, indicating that highly cited papers tend to be cited across different research fields.
}
